\section{Threat Model}
\label{sec:threat}

A claimed irrigation log is a finite map
\begin{equation}
L=\{(t_i,u_i)\}_{i=1}^{|L|},\qquad u_i>0,\ t_i\in[0,T),
\label{eq:log}
\end{equation}
where $u_i$ is a claimed depth in millimetres and $t_i$ an hour index. Three
independent quantities can be forged, and each defence tier observes a different
subset of them:
\begin{equation}
\underbrace{\textstyle\sum_i u_i}_{\text{total volume}}\,,\qquad
\underbrace{\{u_i\}}_{\text{individual depths}}\,,\qquad
\underbrace{\{t_i\}}_{\text{timestamps}}\,.
\label{eq:three}
\end{equation}
Table~\ref{tab:threat} lists the resulting fraud classes. The adversary is
\emph{rational}: it maximises creditable dry phases subject to not being flagged,
and it knows the detector, because the detector is published. Publishing a
threshold is publishing the constraint the adversary optimises against; this is
the sense in which $\varepsilon$ in \eqref{eq:v4} below is not a secret.

Let $M$ be the volume actually pumped, measured at the cooperative station, and
$U(L)=\sum_i u_i$ the claimed total. The tier-1 statistic is
\begin{equation}
G=M-U(L),
\label{eq:gap}
\end{equation}
and the credit-relevant quantity is the number of maximal intervals on which the
replayed depletion stays above $d_{\mathrm{dry}}$ with no claim, of length at
least $72$\,h:
\begin{equation}
n_{\mathrm{dry}}(L)=\bigl|\mathcal{A}\bigl(\mathcal{R}(L)\bigr)\bigr|.
\label{eq:ndry}
\end{equation}
Revenue is proportional to $n_{\mathrm{dry}}$, not to $U(L)$. This distinction is
what makes the attack of Sec.~\ref{sec:timeshift} profitable, and it is invisible
in a formulation that treats total volume as the target.

\section{Physics-Based Filters and Their Exact Blind Spots}
\label{sec:filters}

Replaying $L$ through the FAO-56 root-zone bucket~\cite{fao56} under public ERA5
forcing produces a trajectory $(\hat w_t,\hat d_t)$; the companion
work~\cite{twingate2026} gives the balance in full. Three filters apply.

\emph{$V_1$, saturation.} A claim issued while the root zone is already at field
capacity cannot be stored:
\begin{equation}
V_1(L)=\{\,i:\ u_i>0\ \wedge\ \hat w_{t_i}^{-}\ \ge\ FC-\varepsilon\,\}.
\label{eq:v1}
\end{equation}

\emph{$V_4$, storage excess.} A claim larger than the deficit available at that
hour must have run off or percolated, so it was not ``applied'' in the sense the
standard requires:
\begin{equation}
V_4(L)=\{\,i:\ u_i\ >\ \underbrace{(FC-\hat w_{t_i}^{-})}_{\text{deficit}}
\ +\ \varepsilon\,\}.
\label{eq:v4}
\end{equation}

\emph{$V_5$, infiltration bound.} Independent of state, a single event of
duration $\tau_i$ cannot deliver more than the soil can accept:
\begin{equation}
u_i\ \le\ k_{\mathrm{inf}}\,\tau_i+p_{\max},
\label{eq:v5}
\end{equation}
with $k_{\mathrm{inf}}=10$\,mm\,h$^{-1}$ and $p_{\max}=100$\,mm taken at the
loose end of the published range for puddled rice, so that a farmer who records
two adjacent events in one hour is not accused. At $\tau_i=1$\,h the bound is
$110$\,mm.

A fourth filter, on depletion below the wilting point, was included in the
published version of this layer. It cannot activate: see
Sec.~\ref{sec:v2reach}.

\begin{proposition}[Exact blindness of aggregate defences]
\label{prop:blind}
Let $L$ and $L'$ share the same multiset of depths and the same total, i.e.\
$\{u_i\}=\{u'_j\}$ as multisets. Then
\begin{equation}
G(L)=G(L')\qquad\text{and}\qquad |V_5(L)|=|V_5(L')|.
\end{equation}
\end{proposition}

\begin{proof}
$U(L)=\sum_i u_i$ is a symmetric function of the multiset, so it is invariant
under any reassignment of depths to hours; $M$ is a property of the field, not of
the log; hence $G=M-U$ is unchanged. Likewise \eqref{eq:v5} is evaluated per
event and depends only on $(u_i,\tau_i)$, never on $t_i$ or on $\hat w$, so the
count of violations is a function of the multiset alone.
\end{proof}

Proposition~\ref{prop:blind} is the whole difficulty in one line. The pump meter
and the infiltration bound are the two defences that do not depend on the model
being right, and both are \emph{exactly} invariant under the transformation the
adversary needs. What remains is $V_1$ and $V_4$, which do depend on order
through $\hat w_{t_i}^{-}$. They turn out not to help.

\section{The Time-Shifting Attack}
\label{sec:timeshift}

\begin{proposition}[Credit gain and evasion are the same action]
\label{prop:coincide}
Let a claim at hour $t$ have available deficit $c_t=FC-\hat w_t^{-}$. Between
consecutive claims the bucket loses
$\Delta_t=\sum_{s}(\mathit{ETc}_s+\mathit{DP}_s-P_s^{+})\ge 0$, so delaying a
claim by $\delta$ hours increases its deficit by the drainage accrued in those
hours:
\begin{equation}
c_{t+\delta}\ \ge\ c_t+\min\!\bigl(\delta\cdot\overline{\Delta},\ FC-\hat w_t^{-}\bigr).
\label{eq:delay}
\end{equation}
Consequently the transformation that moves claims \emph{out of} a dry interval
and onto its boundary simultaneously
\begin{enumerate}
\item[(i)] increases $n_{\mathrm{dry}}$, because the interval is no longer
interrupted; and
\item[(ii)] increases $c_t$ at every moved claim, which by \eqref{eq:v4} makes
$V_4$ \emph{less} likely to fire.
\end{enumerate}
There is no trade-off for the adversary to negotiate: the gradient of profit and
the gradient of undetectability point the same way.
\end{proposition}

\begin{proof}
(i) is definitional: $\mathcal{A}$ in \eqref{eq:ndry} requires no claim inside
the interval, so removing one extends or creates a qualifying phase. For (ii),
delaying a claim to $t+\delta$ means the bucket continues to lose water for
$\delta$ further hours with no inflow, so $\hat w_{t+\delta}^{-}\le\hat w_t^{-}$
and $c_{t+\delta}\ge c_t$; the bound \eqref{eq:delay} follows because the loss
per hour is at least the mean $\overline{\Delta}$ until field capacity is
reached, after which $c$ saturates at $FC-\hat w$. Since $V_4$ fires only when
$u_i>c_{t_i}+\varepsilon$ and $u_i$ is unchanged, a larger $c_{t_i}$ can only
reduce the violation set.
\end{proof}

This is why the attack cannot be closed by better hydrology. Any detector whose
verdict is a function of the replayed trajectory must integrate the same
drainage that makes the delayed claim look legitimate. The empirical version of
the attack is a hill-climb over integer hour offsets in $[-36,36]$ that accepts a
move only if the total is preserved bit-for-bit, every depth stays under
\eqref{eq:v5}, and $n_{\mathrm{dry}}$ strictly increases.

\begin{remark}
The attack does not steal water. Total pumped volume is unchanged, so the farmer
pays the same and the atmosphere sees the same methane. What is stolen is
\emph{history}: the credit is issued for a dry-phase pattern that did not occur.
Under VM0051 the crediting unit is precisely that pattern, so the loss is real
even though no physical resource moved.
\end{remark}
